% =====================================================================
%  Trajectory Scheduling:
%  Residue-Driven Process Scheduling and Orthogonal
%  Content--Time Process Annotation
%
%  Self-contained. Every term is defined in the text. Only the
%  standard external literature is cited; no companion paper is
%  presupposed.
% =====================================================================
\documentclass[twocolumn,10pt,a4paper]{article}

\usepackage[utf8]{inputenc}
\usepackage[T1]{fontenc}
\usepackage{lmodern}
\usepackage[a4paper,margin=1.9cm,columnsep=0.6cm]{geometry}
\usepackage{amsmath,amssymb,amsthm,mathtools}
\usepackage{bm}
\usepackage{mathrsfs}
\usepackage{booktabs}
\usepackage{array}
\usepackage{enumitem}
\usepackage{microtype}
\usepackage{algorithm}
\usepackage{algpseudocode}
\usepackage[round,sort&compress]{natbib}
\usepackage[dvipsnames,table]{xcolor}
\usepackage[colorlinks=true,linkcolor=blue!45!black,
            citecolor=blue!45!black,urlcolor=blue!45!black]{hyperref}
\usepackage{cleveref}

% ---- Theorem environments ----
\newtheoremstyle{thmstyle}{\topsep}{\topsep}{\itshape}{}{\bfseries}{.}{.5em}{}
\theoremstyle{thmstyle}
\newtheorem{theorem}{Theorem}[section]
\newtheorem{lemma}[theorem]{Lemma}
\newtheorem{proposition}[theorem]{Proposition}
\newtheorem{corollary}[theorem]{Corollary}
\theoremstyle{definition}
\newtheorem{definition}[theorem]{Definition}
\newtheorem{axiom}{Axiom}
\newtheorem{principle}[theorem]{Principle}
\newtheorem{construction}[theorem]{Construction}
\newtheorem{example}[theorem]{Example}
\theoremstyle{remark}
\newtheorem{remark}[theorem]{Remark}

\crefname{axiom}{Axiom}{Axioms}
\Crefname{axiom}{Axiom}{Axioms}
\crefname{principle}{Principle}{Principles}
\Crefname{principle}{Principle}{Principles}

% ---- Shorthand ----
\newcommand{\R}{\mathbb{R}}
\newcommand{\N}{\mathbb{N}}
\newcommand{\Z}{\mathbb{Z}}
\newcommand{\Zp}{\mathbb{Z}_{\ge 0}}
\newcommand{\floor}{\beta}
\newcommand{\res}{\rho}
\newcommand{\Mcount}{M}
\newcommand{\Tasks}{\mathcal{T}}
\newcommand{\task}{\tau}
\newcommand{\state}{\sigma}
\newcommand{\States}{\Sigma}
\newcommand{\Cls}{\mathsf{C}}
\newcommand{\Addr}{\mathsf{a}}
\newcommand{\dig}{\mathsf{h}}
\newcommand{\lab}{\ell}
\newcommand{\abs}[1]{\lvert#1\rvert}
\newcommand{\norm}[1]{\lVert#1\rVert}
\newcommand{\eps}{\varepsilon}
\newcommand{\dd}{\,\mathrm{d}}
\DeclareMathOperator{\lcp}{lcp}
\DeclareMathOperator{\dec}{dec}
\DeclareMathOperator*{\argmin}{arg\,min}
\DeclareMathOperator*{\argmax}{arg\,max}

\title{\textbf{Trajectory Scheduling:\\[2pt]
Residue-Driven Process Scheduling and Orthogonal\\
Content--Time Process Annotation}}

\author{%
Kundai Farai Sachikonye\\
\small Technical University of Munich\\
\small \texttt{kundai.sachikonye@wzw.tum.de}}

\date{\today}

\begin{document}
\maketitle

% =====================================================================
\begin{abstract}
\noindent
We develop, from first principles, a scheduler that orders work not by
predicted running time---a quantity we prove no scheduler can know in
advance without performing the work---but by a measured \emph{residue}:
the distance of a partially computed result from completion. The
framework rests on a small core of theorems. First, we prove a
\emph{Residue Floor Theorem}: any finite agent that resolves a problem
to a recognisable answer cannot drive its residual error to zero, and a
strictly positive floor $\floor>0$ separates ``recognised as solved''
from ``not yet distinguishable from any other candidate.'' Second, we
prove that the count of committed work units is strictly monotone---no
operation decrements it---so that the natural progress coordinate of a
process is a non-decreasing integer we call its \emph{trajectory count}.
Third, we define a \emph{categorical complexity} that measures a task by
the number of distinguishable work-trajectories required to settle it,
prove a closed-form $\mathcal{O}(\log K)$ termination bound for
content-bounded tasks, and use the live deviation of measured residue
from this baseline to diagnose, at runtime, whether a task is
compute-limited or structure-limited---a distinction we show is the
correct basis for resource allocation. We assemble these into a
residue-driven scheduling algorithm, prove its priority rule sound (it
never starves a task that is converging and never feeds a task that has
stalled), prove a sufficiency-stopping theorem (a sub-task may be
released above its own floor once it crosses the floor of the parent
problem), and bound its overhead. Finally we address process
annotation. We prove a \emph{Label Orthogonality Theorem}: identity and
relatedness are independent requirements served by opposite locality
properties---a cryptographic digest (avalanche: maximal output
divergence from minimal input difference) is optimal for identity and
provably destroys relatedness, while a monotone trajectory address
(prefix-sharing) is optimal for temporal/structural relatedness and
carries no integrity guarantee. We show that a cryptographic hash is, by
construction, time-blind, and that pairing it with the monotone
trajectory address yields a label on which recurrence, co-temporality,
and refinement are each decidable in $\mathcal{O}(1)$, and prove that no
single label can serve both axes. All claims are proved; every
construction is finite; complexity bounds are stated and derived.
\end{abstract}

\noindent\textbf{Keywords:} process scheduling; residue; resolution
floor; anytime computation; categorical complexity; monotone progress;
content addressing; avalanche; locality-sensitive hashing; logical
clocks; sufficiency stopping.

\tableofcontents

% =====================================================================
\section{Introduction}
\label{sec:intro}
% =====================================================================

\subsection{The unknowable-duration problem}

A scheduler decides, repeatedly, which of several pending pieces of work
to advance next, and how much of a shared budget to give it. The
classical theory of scheduling assigns each job a processing time and
optimises a function of completion times \citep{pinedo2016,brucker2007,
liu1973}. The assignment is the difficulty. For a large class of
processes---search, inference, optimisation, any computation whose
running time depends on data it has not yet seen---the processing time is
not known in advance, and cannot be, for a reason that is not
contingent:

\begin{quote}\itshape
If the exact running time of a process were known before running it, the
process would already have been run; the only general way to learn how
long a computation takes is to perform it.
\end{quote}

This is a restatement, in scheduling terms, of the impossibility of a
non-trivial running-time oracle \citep{turing1936,blum1967}. A scheduler
that orders work by predicted duration is therefore ordering by a
quantity it cannot measure, and must fall back on heuristics, historical
averages, or static priorities, each of which is wrong whenever the data
diverges from the assumption.

\subsection{The move: schedule by closeness, not by duration}

We make one change of variable. Instead of asking \emph{how long will
this take}---unknowable---we ask \emph{how close is this to done}---a
property of the present partial result, measurable now. The two are not
the same question, and the second is the better one for a scheduler:
remaining \emph{duration} is a prediction about the future; remaining
\emph{distance to a solution} is a measurement of the present.

To make ``distance to a solution'' a rigorous quantity we need three
things, each of which we derive from scratch:

\begin{enumerate}[leftmargin=1.4em,itemsep=2pt]
\item A notion of how far a partial result stands from being a
\emph{recognised} answer, and a proof that this distance has a strictly
positive floor---that ``solved'' is a region, not a point
(\Cref{sec:floor}).
\item A monotone progress coordinate: a count of committed work that
never decreases, so that ``how far along'' is well-defined and a
revisited state is never a return (\Cref{sec:count}).
\item A complexity measure, computed before running, that predicts how
many work units a task \emph{should} need, against which the live
measurement is compared to diagnose the kind of limit the task faces
(\Cref{sec:complexity}).
\end{enumerate}

\subsection{Two products}

The paper has two deliverables. The first is a \emph{residue-driven
scheduler} (\Cref{sec:scheduler}): an algorithm that prioritises tasks
by measured residue and its rate of descent, allocates compute to tasks
whose residue is falling (compute-limited) and withholds it from tasks
whose residue has gone flat (structure-limited), and stops a sub-task as
soon as it is good enough for the larger problem it serves, even if it
has not reached its own floor. We prove the priority rule sound and the
sufficiency-stop correct.

The second is a \emph{process annotation scheme}
(\Cref{sec:annotation}). A scheduler must label processes so as to answer
two questions in $\mathcal{O}(1)$: \emph{is this the same work as that}
(identity), and \emph{is this related to that} (relatedness). We prove
these require opposite locality and therefore cannot be served by one
label: a cryptographic digest, whose avalanche property maps minimal
input differences to maximal output differences \citep{shannon1949,
webster1986,feistel1973}, is optimal for identity and provably destroys
relatedness; a monotone trajectory address, whose prefixes are shared by
related work, is optimal for relatedness and carries no integrity. We
prove further that a cryptographic hash is structurally time-blind---time
is not among its inputs---and that pairing it with the monotone progress
coordinate of \Cref{sec:count} supplies exactly the missing axis, making
recurrence, co-temporality, and refinement decidable by direct
comparison.

\subsection{What is and is not claimed}

We do not claim to predict running time; we prove it cannot be done and
build a scheduler that does not need it. We do not claim a new
cryptographic primitive; we use standard hashes as black boxes and prove
a separation between what they can and cannot annotate. We do not claim
to reduce any standard complexity class; the categorical complexity of
\Cref{sec:complexity} is a refinement along an access-pattern axis, with
explicit non-tight relations to time and space
(\Cref{sec:complexity-standard}). Every theorem is proved from the stated
definitions, and where a claim could not be proved at this standard it is
omitted.

\subsection{Organisation}

\Cref{sec:model} fixes the agent, the candidate space, and the residue
functional. \Cref{sec:floor} proves the Residue Floor Theorem.
\Cref{sec:count} proves strict monotonicity of the trajectory count.
\Cref{sec:complexity} develops categorical complexity and the
$\mathcal{O}(\log K)$ termination bound. \Cref{sec:diagnosis} proves the
compute-versus-structure diagnosis. \Cref{sec:scheduler} states and
analyses the scheduler. \Cref{sec:sufficiency} proves sufficiency
stopping. \Cref{sec:annotation} proves the Label Orthogonality and
Time-Blindness theorems and constructs the paired annotation.
\Cref{sec:overhead} bounds total overhead. \Cref{sec:related} situates
the work. \Cref{sec:conclusion} concludes.

% =====================================================================
\section{The model}
\label{sec:model}
% =====================================================================

We fix a domain-agnostic model of a computation that resolves a problem
to an answer. Nothing in it is specific to any machine.

\subsection{Candidates, truth, and the resolving agent}

\begin{definition}[Problem instance]
\label{def:instance}
A \emph{problem instance} is a pair $(\mathcal{X}, \mathcal{X}^\ast)$
where $\mathcal{X}$ is a non-empty set of \emph{candidates} (admissible
partial or complete answers) and $\varnothing\neq\mathcal{X}^\ast
\subseteq\mathcal{X}$ is the set of \emph{solutions}. We allow
$\abs{\mathcal{X}^\ast}>1$: several candidates may be operationally
correct.
\end{definition}

\begin{definition}[Finite resolving agent]
\label{def:agent}
A \emph{finite resolving agent} $A$ is a tuple
$(K_A,\,\Phi_A,\,\dec_A)$ where $K_A$ is a finite set of
\emph{internal states} (the agent's representational capacity),
$\Phi_A\colon\mathcal{X}\to K_A$ encodes a candidate into a state, and
$\dec_A\colon K_A\to\mathcal{X}$ decodes a state back to a candidate. The
agent is \emph{bounded} for the instance if $\abs{K_A}<\abs{\mathcal{X}}$
(strictly fewer internal states than candidates).
\end{definition}

Boundedness is the only substantive hypothesis on the agent. It holds for
every physically realised computer relative to any candidate space large
enough to matter (finite memory, unbounded or merely larger candidate
space).

\subsection{The residue functional}

The agent measures how far a candidate stands from solving the instance.

\begin{definition}[Residue functional]
\label{def:residue}
A \emph{residue functional} for agent $A$ is a map
\[
\res_A\colon \mathcal{X}\times\mathcal{X}^\ast \to [0,U],
\qquad U>0,
\]
where $\res_A(x,x^\ast)$ is the agent's measured misalignment of
candidate $x$ from solution $x^\ast$, satisfying:
\begin{enumerate}[label=\textup{(R\arabic*)},leftmargin=2.8em,itemsep=2pt]
\item\label{R:nonneg} \textbf{Non-negativity.} $\res_A(x,x^\ast)\ge 0$.
\item\label{R:bound} \textbf{Boundedness.} $\res_A(x,x^\ast)\le U$.
\item\label{R:coherent} \textbf{Representational coherence.} There is a
function $\delta_A\colon K_A\times K_A\to[0,U]$ with $\delta_A(k,k)=\floor_A$
constant on the diagonal, such that
$\res_A(x,x^\ast)=\delta_A(\Phi_A(x),\Phi_A(x^\ast))$. That is, the agent
reads residue only through its internal encodings of the two candidates,
and two candidates it encodes identically have the diagonal value.
\end{enumerate}
\end{definition}

\ref{R:coherent} is the operational content: an agent cannot compare
candidates except through its own finite representation of them. The
constant $\floor_A:=\delta_A(k,k)$ is the residue the agent reads between
a state and itself; we will prove it is strictly positive and that it is
the floor of $\res_A$.

\begin{remark}[Why residue, not ``time remaining'']
\label{rem:why-residue}
$\res_A(x,x^\ast)$ depends only on the present candidate $x$ and is
evaluable now; ``time remaining'' depends on the future trajectory and is
not. The entire design rests on substituting the first for the second.
\end{remark}

% =====================================================================
\section{The Residue Floor Theorem}
\label{sec:floor}
% =====================================================================

We prove that a bounded agent cannot reduce residue to zero, and that the
floor $\floor_A>0$ is the boundary between ``recognised as a solution''
and ``indistinguishable from any candidate.'' This is the theorem that
makes ``solved'' a region rather than a point, and hence makes residue a
usable scheduling signal.

\begin{theorem}[Residue Floor]
\label{thm:floor}
For every bounded resolving agent $A$ with residue functional $\res_A$,
the diagonal value $\floor_A$ is strictly positive, and no candidate
attains residue below it:
\[
\inf_{x\in\mathcal{X}}\res_A(x,x^\ast)=\floor_A>0
\qquad\text{for every }x^\ast\in\mathcal{X}^\ast .
\]
\end{theorem}

\begin{proof}
Suppose $\floor_A=0$. Then by \ref{R:coherent} there is, for some
$x^\ast$, a candidate $x_0$ with $\res_A(x_0,x^\ast)=0$, hence
$\delta_A(\Phi_A(x_0),\Phi_A(x^\ast))=0$. Since the diagonal value is
$\floor_A=0$ and (by \ref{R:coherent}) the diagonal is where $\delta_A$
attains its minimum reading, $\delta_A(\Phi_A(x_0),\Phi_A(x^\ast))=0$
forces $\Phi_A(x_0)=\Phi_A(x^\ast)$: the agent encodes $x_0$ and
$x^\ast$ to the \emph{same} internal state.

Now count. The agent's representable candidates are
$\dec_A(\Phi_A(\mathcal{X}))$, of cardinality at most $\abs{K_A}$. Because
$A$ is bounded, $\abs{K_A}<\abs{\mathcal{X}}$, so $\Phi_A$ is not
injective and some distinct candidates share an internal state. Zero
residue, by the above, means the agent cannot internally separate $x_0$
from $x^\ast$; but then it equally cannot separate $x_0$ from \emph{any}
candidate sharing that state. The reading ``$x_0$ is the solution'' is
then indistinguishable, within $A$, from ``$x_0$ is any of the candidates
collapsed onto $\Phi_A(x^\ast)$.'' A residue of $0$ thus carries no
information identifying $x_0$ as the solution rather than as an
arbitrary member of its collision class. This contradicts the premise
that $\res_A=0$ marks recognition of the solution: a value that does not
distinguish the solution from non-solutions does not recognise it. Hence
$\floor_A>0$. That $\floor_A$ is the infimum is \ref{R:coherent}: every
reading factors through $\delta_A$, whose minimum is the diagonal value
$\floor_A$.
\end{proof}

\begin{corollary}[Recognition requires a gap]
\label{cor:gap}
A candidate is recognisable as a solution by $A$ only at residue
$\ge\floor_A>0$. The positive gap is not error to be eliminated; it is
the margin by which the solution stands apart from the candidates the
agent cannot distinguish from it. Driving residue to zero destroys
recognition rather than perfecting it.
\end{corollary}

\begin{proof}
Immediate from \Cref{thm:floor}: below $\floor_A$ the solution collapses,
in the agent's representation, into a class it cannot resolve.
\end{proof}

\begin{theorem}[Information capacity of the residue scale]
\label{thm:info}
At reading precision $\eps>0$, the number of distinguishable residue
values an agent can record is at most $\lceil(U-\floor_A)/\eps\rceil$,
carrying at most $\log_2\!\big((U-\floor_A)/\eps\big)$ bits.
\end{theorem}

\begin{proof}
By \Cref{thm:floor} the image of $\res_A$ lies in $[\floor_A,U]$, an
interval of length $U-\floor_A$; partitioning into bins of width $\eps$
gives the count, and the binary logarithm gives the bit bound
\citep{shannon1948,cover2006}.
\end{proof}

\Cref{thm:info} is a sizing law used later: a scheduler that ranks tasks
by residue can distinguish at most $\mathcal{O}((U-\floor_A)/\eps)$
priority levels, so its comparisons are over a finite, bounded-precision
ladder, and ties are genuine (two tasks may be residue-indistinguishable).

% =====================================================================
\section{The trajectory count is strictly monotone}
\label{sec:count}
% =====================================================================

We now fix the progress coordinate. A process advances by committing
discrete \emph{work units}; we prove the count of committed units is
strictly increasing, so it is a faithful, non-decreasing clock against
which residue descent is measured, and a revisited candidate is never a
return to a prior process state.

\subsection{Work units and the count}

\begin{definition}[Work unit; trajectory count]
\label{def:count}
A \emph{work unit} of a process is a committed step that changes the
process's record: it appends one entry to an append-only log of the
process's history. The \emph{trajectory count} $\Mcount(t)\in\Zp$ is the
number of work units committed up to step $t$. The process \emph{state}
at step $t$ is the pair
\[
s_t=\big(x_t,\ \Mcount(t)\big),
\]
the current candidate together with the count; two states are equal iff
both components agree.
\end{definition}

\begin{theorem}[Strict monotonicity]
\label{thm:monotone}
Along any process the trajectory count is strictly increasing: each
committed work unit increments $\Mcount$ by exactly one, so for $i<j$,
$\Mcount(i)<\Mcount(j)$. No operation decrements $\Mcount$.
\end{theorem}

\begin{proof}
By \Cref{def:count} each work unit appends exactly one log entry and
$\Mcount$ counts entries; appending increases the count by one. The log
is append-only, so no entry is removed and the count never decreases. A
purported ``undo'' is itself a committed step---it appends a
compensating entry---hence increments $\Mcount$ rather than decrementing
it: undoing is a forward unit, not a reversal of the count.
\end{proof}

\begin{corollary}[No return; resemblance is not identity]
\label{cor:noreturn}
For $i<j$, $s_i\neq s_j$ even when $x_i=x_j$: a process that revisits a
candidate is in a new state, because its count has strictly grown. Two
process states with identical candidates but different counts are
distinct.
\end{corollary}

\begin{proof}
By \Cref{thm:monotone}, $\Mcount(i)<\Mcount(j)$, so the second components
of $s_i,s_j$ differ; hence the states differ regardless of the
candidates.
\end{proof}

\begin{remark}[Why a count, not a wall clock]
\label{rem:why-count}
The trajectory count orders a process's own work intrinsically. Unlike a
wall clock it cannot run backward, cannot be reset, and cannot collide
(two units in the same physical instant still receive distinct, ordered
counts). It is the process's logical time
\citep{lamport1978,fidge1988,mattern1989}, and---being monotone by
construction (\Cref{thm:monotone})---it is the coordinate against which a
\emph{decreasing} residue is meaningful: residue measured per unit count
is the descent rate driving the scheduler (\Cref{sec:scheduler}).
\end{remark}

% =====================================================================
\section{Categorical complexity and the termination bound}
\label{sec:complexity}
% =====================================================================

Residue tells the scheduler how far a task is from done \emph{now};
categorical complexity tells it how far the task \emph{should} have to
travel, as a static prediction, so that the deviation between the two can
be read at runtime. We define complexity by the number of distinguishable
work-trajectories a task requires, derive its closed form, and prove a
logarithmic termination bound for the content-bounded case.

\subsection{Work-trajectories and their count}

\begin{definition}[Work-trajectory; type]
\label{def:trajectory}
A \emph{work-trajectory} of length $n$ is an ordered sequence of $n$
committed work units. Each unit carries one of $d$ \emph{operation types}
from a fixed finite set (e.g.\ read, transform, branch, record). The
\emph{commitment structure} of a trajectory is a partition of its $n$
units into maximal contiguous \emph{blocks} the process records as single
checkpoints; a length-$n$ trajectory with $k$ blocks is recorded as a
composition $(c_1,\dots,c_k)$ of $n$ together with a type for each block.
\end{definition}

\begin{lemma}[Compositions of an integer]
\label{lem:compositions}
The number of ordered compositions of $n$ into $k$ positive parts is
$\binom{n-1}{k-1}$; summed over $k$, the total number of compositions of
$n$ is $2^{n-1}$ \citep{stanley2012}.
\end{lemma}

\begin{proof}
A composition into $k$ parts is a choice of $k-1$ of the $n-1$ gaps
between $n$ ordered units as block boundaries, giving $\binom{n-1}{k-1}$;
$\sum_{k=1}^{n}\binom{n-1}{k-1}=\sum_{j=0}^{n-1}\binom{n-1}{j}=2^{n-1}$
by the binomial theorem.
\end{proof}

\begin{theorem}[Trajectory inflation]
\label{thm:inflation}
The number of distinguishable work-trajectories of length $n$ over $d$
operation types, counted with their commitment structure, is
\[
\boxed{\,T(n,d)=d\,(d+1)^{\,n-1}.\,}
\]
\end{theorem}

\begin{proof}
A distinguishable trajectory is a composition of $n$ into $k$ blocks
(\Cref{lem:compositions}) together with a choice of one of $d$ types per
block, giving $d^k$ labellings for a $k$-block composition. Summing over
commitment depth $k$,
\[
T(n,d)=\sum_{k=1}^{n}\binom{n-1}{k-1}d^{k}
= d\sum_{j=0}^{n-1}\binom{n-1}{j}d^{j}
= d\,(1+d)^{n-1},
\]
the inner sum being $(1+d)^{n-1}$ by the binomial theorem.
\end{proof}

\begin{corollary}[Multiplicative recurrence]
\label{cor:recurrence}
$T(n{+}1,d)=(d+1)\,T(n,d)$ with $T(1,d)=d$: each additional committed
unit multiplies the trajectory count by $d+1$ (the $d$ type choices for a
new block, plus the single option of extending the current block).
\end{corollary}

\begin{proof}
$T(n{+}1,d)=d(d+1)^{n}=(d+1)T(n,d)$; $T(1,d)=d(d+1)^0=d$.
\end{proof}

\subsection{Categorical complexity and termination}

\begin{definition}[Categorical complexity]
\label{def:complexity}
The \emph{categorical complexity} of a task is the number $K\ge 1$ of
distinguishable work-trajectories that must be realised before the task's
residue is certified at the floor---i.e.\ the size of the trajectory
space the task must explore to recognise its solution. A task is
\emph{content-bounded} if its completion condition is ``$K$ trajectories
realised'' (as opposed to ``an external event has occurred,'' treated in
\Cref{sec:diagnosis}).
\end{definition}

\begin{theorem}[Logarithmic termination bound]
\label{thm:termination}
A content-bounded task of categorical complexity $K$ over $d$ operation
types reaches completion after committing
\[
n_{\mathrm{term}}(K,d)=1+\Big\lceil\log_{d+1}\!\tfrac{K}{d}\Big\rceil
\]
work units. Its cumulative committed cost is at least
$n_{\mathrm{term}}\cdot\floor_A=\Theta(\log K)$.
\end{theorem}

\begin{proof}
By \Cref{thm:inflation}, $n$ committed units realise $T(n,d)=d(d+1)^{n-1}$
distinguishable trajectories. The task completes at the least $n$ with
$T(n,d)\ge K$, i.e.\ $d(d+1)^{n-1}\ge K$, i.e.\ $(d+1)^{n-1}\ge K/d$,
i.e.\ $n\ge 1+\log_{d+1}(K/d)$; the least integer is
$n_{\mathrm{term}}=1+\lceil\log_{d+1}(K/d)\rceil$. Each committed unit
deposits residue at least $\floor_A$ (\Cref{thm:floor} read as a
per-unit lower bound on resolved content), so cumulative cost is
$\ge n_{\mathrm{term}}\floor_A=\Theta(\log K)$.
\end{proof}

\begin{remark}[Linear cost, exponential content]
\label{rem:asymmetry}
The cumulative \emph{cost} of $n$ units grows linearly ($\ge n\floor_A$),
while the \emph{content}---distinguishable trajectories---grows
exponentially ($T(n,d)=d(d+1)^{n-1}$). A task that terminates on content
(``$K$ trajectories'') terminates in $\mathcal{O}(\log K)$ units; a task
that terminates on cost (``deposit total residue $R$'') terminates in
$\mathcal{O}(R)$ units. The two are different stopping rules, and a task's
specification, not the framework, selects which governs it
(\Cref{sec:diagnosis}).
\end{remark}

\subsection{Relation to standard complexity}
\label{sec:complexity-standard}

\begin{proposition}[Non-tight relations]
\label{prop:standard}
Categorical complexity refines time/space complexity along an
access-pattern axis, with non-tight inclusions. A task content-bounded
with $K$ polynomial in instance size is solvable in time polynomial in
$\log K$ and hence polynomial in input size; a task with no
trajectory structure (every candidate must be enumerated) has
$n_{\mathrm{term}}=\Theta(K)$, matching brute force. Categorical
simplicity neither implies nor is implied by polynomial-time simplicity.
\end{proposition}

\begin{proof}
The first claim is \Cref{thm:termination} with $K=\mathrm{poly}$:
$n_{\mathrm{term}}=\mathcal{O}(\log K)$ units, each of bounded cost. The
second: if no proper subset of trajectories certifies the solution, the
completion condition is ``all $K$ realised,'' so $T(n,d)\ge K$ is needed
with the trajectories enumerated, giving $\Theta(K)$. Independence: a
problem can be categorically simple (one trajectory settles it) yet
costly per unit, or categorically hard yet cheap per unit; the two axes
measure different resources \citep{hartmanis1965,arora2009,sipser2013}.
No claim of class collapse is made or used.
\end{proof}

% =====================================================================
\section{Runtime diagnosis: compute-limited vs.\ structure-limited}
\label{sec:diagnosis}
% =====================================================================

The static complexity $K$ is a prediction; the live residue is a
measurement. Their relationship at runtime classifies the kind of limit a
task faces, which we prove is the correct basis for allocation.

\subsection{Descent and the predicted curve}

\begin{definition}[Residue descent; on-curve]
\label{def:descent}
For a running task, write $\res(n)$ for its residue after $n$ committed
units. Its \emph{descent rate} at $n$ is
$\Delta(n)=\res(n-1)-\res(n)\ge 0$ (residue is non-increasing along a
correct process: committing work cannot increase distance to a solution
the process is converging to). The task is \emph{on-curve} if its
realised descent matches the rate implied by its categorical complexity
$K$ (\Cref{thm:termination}): residue should fall from its start toward
$\floor_A$ over $n_{\mathrm{term}}(K,d)$ units. It is \emph{behind curve}
if $\res(n)$ exceeds the predicted residue at the same $n$.
\end{definition}

\begin{definition}[Limit type]
\label{def:limit}
A running task above its floor ($\res(n)>\floor_A$) is:
\begin{itemize}[leftmargin=1.4em,itemsep=2pt]
\item \emph{compute-limited} if $\Delta(n)>0$: it is still converging;
the only obstacle to completion is further committed work.
\item \emph{structure-limited} if $\Delta(n)=0$ over a window of $w$
units (a \emph{stall}): residue has stopped falling; further compute does
not reduce distance to a solution.
\end{itemize}
\end{definition}

\begin{theorem}[Diagnosis correctness]
\label{thm:diagnosis}
A task's limit type is decidable from its residue stream alone, without
knowing its running time, and the decision is correct in the following
sense:
\begin{enumerate}[label=\textup{(\roman*)},leftmargin=2.2em,itemsep=2pt]
\item if $\Delta(n)>0$ then allocating one more unit strictly reduces
residue, so additional compute makes progress toward the floor;
\item if $\Delta(n)=0$ over a window then no allocation of further
units within the task's current structure reduces residue, so additional
compute makes no progress.
\end{enumerate}
\end{theorem}

\begin{proof}
(i) $\Delta(n)>0$ is, by \Cref{def:descent}, the statement that the most
recent unit strictly reduced residue; by the process's determinism the
next unit of the same kind reduces it again unless the task crosses the
floor (\Cref{thm:floor}) or stalls---both of which are detected as
$\Delta=0$ at the next step. Hence while $\Delta>0$, compute is exactly
the operative limit.
(ii) $\Delta(n)=0$ over a window means the committed units in that window
left residue unchanged: the candidate moved (count advanced,
\Cref{thm:monotone}) but did not approach a solution. By
\Cref{def:residue} residue is a function of the candidate's encoding; an
unchanged residue under changing candidates means the task is cycling
within a residue level-set, which further same-structure units cannot
leave. Compute is therefore not the limit; the task is bounded by its
structure (missing input, wrong decomposition, or genuine
non-convergence). Both readings use only $\res(\cdot)$, which is
measurable at each unit, and never the running time.
\end{proof}

\begin{remark}[Behind-curve is reclassification, not failure]
\label{rem:reclassify}
A task with $n>n_{\mathrm{term}}(K,d)$ that has not completed is
\emph{not} ipso facto failed: by \Cref{prop:standard} the static $K$ may
have under-estimated the task's true categorical class. The correct
response is to re-estimate $K$ from the realised descent and re-apply
\Cref{thm:diagnosis}: if still descending, the task is compute-limited
and merely harder than predicted (allocate more); only if stalled is it
structure-limited (decline). Treating prediction-overrun as automatic
failure---a common scheduler heuristic---discards converging work.
\end{remark}

% =====================================================================
\section{The residue-driven scheduler}
\label{sec:scheduler}
% =====================================================================

We now assemble the scheduler. It maintains a set of live tasks, each
with a measured residue stream and a static complexity; on each tick it
selects tasks by a residue-based priority, allocates a share of the
tick's budget, dispatches a work unit, reads the returned residue, and
repeats.

\subsection{State and priority}

\begin{definition}[Live task]
\label{def:livetask}
A \emph{live task} $\task$ carries: a categorical complexity $K_\task$
(with $d$ fixed); a committed-unit count $n_\task$ (its trajectory count,
\Cref{def:count}); a current residue $\res_\task=\res_\task(n_\task)$; a
recent descent rate $\Delta_\task$ over a fixed window; a sufficiency
threshold $\theta_\task\in[\floor_A,U]$ (\Cref{sec:sufficiency}); and a
state in $\{\textsc{running},\textsc{stalled},\textsc{sufficient},
\textsc{done},\textsc{declined}\}$.
\end{definition}

\begin{definition}[Residue priority]
\label{def:priority}
The \emph{priority} of a running task $\task$ is
\[
P(\task)=
\frac{\Delta_\task}{\max(\res_\task-\theta_\task,\ \floor_A)} .
\]
If $\res_\task\le\theta_\task$ (the task has reached its sufficiency
threshold), $P(\task)=+\infty$ (dispatch to finalise immediately). If
$\Delta_\task=0$ over the stall window, $P(\task)=0$ (do not feed a
stalled task).
\end{definition}

The numerator rewards \emph{descent rate} (a fast-converging task is
worth advancing); the denominator rewards \emph{closeness} to the
sufficiency threshold (a nearly-done task is worth finishing). A stalled
task ($\Delta=0$) gets priority $0$ regardless of closeness---compute
will not help it. A task at threshold gets $+\infty$---finish it and free
its resources. This realises the two scheduling instincts proved correct
in \Cref{thm:diagnosis}: feed the converging, finish the close, starve the
stalled.

\begin{remark}[Why not $1/n_{\mathrm{term}}$]
\label{rem:not-static}
A scheduler could rank by static remaining units
$n_{\mathrm{term}}(K)-n_\task$. This is a prediction and inherits the
unknowable-duration problem: a behind-curve task (\Cref{rem:reclassify})
would be deprioritised exactly when it is converging and should be fed.
\Cref{def:priority} uses the \emph{measured} $\Delta$ and $\res$, so it is
a present-tense quantity, not a forecast.
\end{remark}

\subsection{The tick loop}

\begin{algorithm}[H]
\caption{Residue-driven scheduler tick}
\label{alg:tick}
\begin{algorithmic}[1]
\Function{Tick}{$L,\,B$} \Comment{live tasks $L$, tick budget $B$}
  \State $R\gets\varnothing$ \Comment{collected results}
  \While{$B>0$ \textbf{and} $\exists\,\task\in L:\task.\text{state}=\textsc{running}$}
    \State $\task\gets\argmax_{\task'\in L}P(\task')$
    \If{$P(\task)=0$} \Comment{all runnable tasks stalled}
      \State \textbf{break}
    \EndIf
    \State $s\gets P(\task)/\textstyle\sum_{\task'}\max(P(\task'),0)$
    \State $b\gets\min(B,\ \lceil s\cdot B\rceil)$ \Comment{budget share}
    \State $r\gets \textsc{Dispatch}(\task,\,b)$ \Comment{one work unit}
    \State $\task.n\mathrel{+}=1$;\ \ $\task.\res\gets r.\res$
    \State update $\task.\Delta$ over the window
    \State append $r$ to the append-only log;\ \ $R\gets R\cup\{r\}$
    \State $B\mathrel{-}=b$
    \If{$\task.\res\le\theta_\task$} $\task.\text{state}\gets\textsc{sufficient}$
    \ElsIf{$\task.\Delta=0$ over window} $\task.\text{state}\gets\textsc{stalled}$
    \EndIf
  \EndWhile
  \State \Return $R$
\EndFunction
\end{algorithmic}
\end{algorithm}

\begin{theorem}[Priority-rule soundness]
\label{thm:sound}
\Cref{alg:tick} (i) never dispatches a stalled task, (ii) never withholds
dispatch from the unique converging task when budget remains, and (iii)
dispatches a task that has reached its threshold before any task that has
not.
\end{theorem}

\begin{proof}
(i) A stalled task has $\Delta=0$, hence $P=0$ by \Cref{def:priority};
\textsc{Tick} selects $\argmax P$ and breaks when the maximum is $0$, so a
stalled task is dispatched only if it is the unique runnable task and even
then the loop breaks before dispatching it. (ii) If exactly one task is
converging ($\Delta>0$) and others are stalled, the converging task is the
unique non-zero priority, so $\argmax P$ selects it while $B>0$. (iii) A
task at threshold has $P=+\infty$, strictly exceeding every finite
priority, so $\argmax P$ selects it first.
\end{proof}

\begin{theorem}[No livelock; progress or decline]
\label{thm:progress}
On each tick, either some task's trajectory count strictly increases, or
every runnable task is stalled and the scheduler signals decline. The
trajectory count of the whole system (sum over tasks) is strictly
monotone across non-declining ticks.
\end{theorem}

\begin{proof}
If any runnable task has $P>0$, \textsc{Tick} dispatches it, incrementing
its count (\Cref{thm:monotone}); the system count rises. If all runnable
tasks have $P=0$ they are stalled; the loop breaks without dispatch and
the scheduler reports decline for the stalled set. Monotonicity of the
system count across dispatching ticks follows from per-task monotonicity
and the append-only log.
\end{proof}

\begin{remark}[Decline as an honest output]
\label{rem:decline}
\Cref{thm:progress} makes ``stuck'' a detectable state---all runnable
tasks stalled---rather than an elapsed timeout. The scheduler distinguishes
``slow but converging'' (keep going) from ``not converging'' (decline)
using the descent signal, never a wall-clock deadline. This is the
anytime-computation stance \citep{dean1988,zilberstein1996} sharpened: the
stopping criterion is non-approach, not time spent.
\end{remark}

% =====================================================================
\section{Sufficiency stopping}
\label{sec:sufficiency}
% =====================================================================

A sub-task serves a larger task. Its job is not to reach its \emph{own}
floor but to be good enough for the \emph{parent}. We formalise this and
prove that a sub-task may be released above its own floor, and that this
judgment cannot be made by the sub-task itself.

\subsection{Parent and sub-task floors}

\begin{definition}[Task graph; parent floor]
\label{def:taskgraph}
A \emph{task graph} is a directed acyclic graph whose vertices are tasks
and whose arcs $\task'\!\to\!\task$ mean ``$\task$ consumes $\task'$'s
output.'' A designated sink carries the global goal $g$ with its own
residue functional $\res_g$ and floor $\floor_g$. For a sub-task $\task$
feeding (directly or transitively) the sink, the \emph{sufficiency
threshold} $\theta_\task$ is the largest residue of $\task$'s output at
which the sink's attainable residue toward $g$ is unaffected:
\[
\theta_\task=\sup\big\{\,r:\ \res_g^{\min}(\task\!=\!r)=\res_g^{\min}(\task\!=\!\floor_A)\big\},
\]
where $\res_g^{\min}(\cdot)$ is the best residue the rest of the graph can
reach toward $g$ given $\task$'s output residue.
\end{definition}

\begin{theorem}[Sufficiency stopping]
\label{thm:sufficiency}
A sub-task $\task$ may be released (marked \textsc{sufficient}) as soon as
$\res_\task\le\theta_\task$, even when $\res_\task>\floor_A$ (it has not
reached its own floor). Releasing at $\theta_\task$ does not change the
global attainable residue toward $g$.
\end{theorem}

\begin{proof}
By \Cref{def:taskgraph}, $\theta_\task$ is exactly the threshold below
which further reduction of $\task$'s residue leaves $\res_g^{\min}$
unchanged: the sink is already at the residue it would reach were
$\task$ resolved to its own floor. Hence committing further units to
$\task$ below $\theta_\task$ reduces $\res_\task$ but not $\res_g^{\min}$;
the marginal contribution to the global goal is zero. Releasing $\task$ at
$\theta_\task$ therefore preserves $\res_g^{\min}$ while saving the
committed cost of driving $\task$ from $\theta_\task$ to $\floor_A$.
\end{proof}

\begin{example}[Coarse-but-sufficient]
\label{ex:timetable}
Let $g$ be ``produce a train timetable,'' whose floor $\floor_g$ is at
minute resolution, and let $\task$ be a timing sub-task capable of
microsecond residue. The sink's residue toward $g$ is unchanged once
$\task$'s output is correct to the minute, so $\theta_\task$ sits at
minute resolution---far above $\task$'s own microsecond floor. By
\Cref{thm:sufficiency}, $\task$ is released at the minute, and the
microsecond precision it could attain is never computed. The extra
precision is residue no part of $g$ registers; computing it is the cost of
closing a gap the parent does not need closed (\Cref{cor:gap}).
\end{example}

\subsection{Sufficiency is not locally decidable}

\begin{theorem}[No local sufficiency]
\label{thm:nolocal}
A sub-task $\task$ cannot, in general, compute its own sufficiency
threshold $\theta_\task$. The threshold depends on $\res_g^{\min}$, a
function of the global goal $g$ and the whole task graph; if $g$ is not
represented within $\task$, no quantity internal to $\task$ is a function
of $\theta_\task$.
\end{theorem}

\begin{proof}
By \Cref{def:taskgraph}, $\theta_\task$ is defined through
$\res_g^{\min}(\cdot)$, which quantifies over the rest of the graph and
the goal $g$. If $g\notin$ representation of $\task$, then $\task$ holds
no value depending on $\res_g$: two task graphs agreeing on $\task$'s
local data but differing in the downstream structure that fixes
$\res_g^{\min}$ assign different $\theta_\task$, yet are indistinguishable
to $\task$. Hence $\theta_\task$ is not a function of $\task$-internal
data, and $\task$ cannot compute it.
\end{proof}

\begin{corollary}[The threshold is set above the sub-task]
\label{cor:above}
Sufficiency thresholds must be assigned by a layer that holds the global
goal and the task graph---the scheduler/orchestrator---and supplied to
each sub-task as data. The sub-task reports its residue; the supervising
layer decides whether that residue is sufficient for $g$ and signals
release.
\end{corollary}

\begin{proof}
Immediate from \Cref{thm:nolocal}: the only observer able to evaluate
$\res_g^{\min}$ is one holding $g$ and the graph, which is not any single
sub-task.
\end{proof}

This is the one place the global goal touches the scheduler's per-task
loop: the threshold $\theta_\task$ (a number) is supplied top-down, while
the residue $\res_\task$ flows bottom-up; \Cref{alg:tick} compares them.
The separation keeps the scheduler's measurement local (it reads only
$\res_\task,\Delta_\task$) and the sufficiency judgment global.

% =====================================================================
\section{Process annotation: identity and relatedness}
\label{sec:annotation}
% =====================================================================

A scheduler must label every process so as to answer, in $\mathcal{O}(1)$,
two questions: \emph{is this the same work as that} (identity) and
\emph{is this related to that} (relatedness, e.g.\ shares structure, or
occurs in the same burst). We prove these demand opposite locality and
therefore cannot be served by one label, characterise the two labels that
serve them, and construct their pairing.

\subsection{Two labelling requirements}

\begin{definition}[Labelling; locality]
\label{def:label}
A \emph{labelling} is a map $\lab\colon\mathcal{P}\to\mathcal{L}$ from
processes (each a content with a trajectory state) to labels, with a
label metric $d_\mathcal{L}$. The labelling is:
\begin{itemize}[leftmargin=1.4em,itemsep=2pt]
\item \emph{identity-faithful} if distinct processes receive labels at
near-maximal $d_\mathcal{L}$ (no two processes are confusable), and
collisions are cryptographically hard to produce;
\item \emph{relatedness-faithful} if processes that are close in the
process metric (shared structure, or temporal proximity) receive labels
close in $d_\mathcal{L}$.
\end{itemize}
\end{definition}

\begin{theorem}[Label orthogonality]
\label{thm:orthogonality}
No single labelling is both identity-faithful and relatedness-faithful.
The two properties impose opposite locality: identity-faithfulness sends
near inputs to far labels; relatedness-faithfulness sends near inputs to
near labels.
\end{theorem}

\begin{proof}
Let $\lab$ be identity-faithful. Take two processes $p,p'$ that are
\emph{related} (close in the process metric) but distinct---e.g.\ two
inputs differing in one bit. Identity-faithfulness requires
$d_\mathcal{L}(\lab(p),\lab(p'))$ near-maximal, since $p\neq p'$ and the
labelling must not confuse distinct processes; this is the avalanche
property \citep{shannon1949,webster1986,feistel1973}: minimal input
difference maps to maximal output difference. But relatedness-faithfulness
requires $d_\mathcal{L}(\lab(p),\lab(p'))$ \emph{small}, because $p,p'$
are close. The two demands on $d_\mathcal{L}(\lab(p),\lab(p'))$---near-%
maximal and small---are contradictory for the same related pair. Hence no
$\lab$ satisfies both.
\end{proof}

\begin{corollary}[A cryptographic digest cannot annotate relatedness]
\label{cor:no-related}
A cryptographic hash $\dig$ (designed for the strict avalanche criterion
\citep{webster1986}, hence identity-faithful) is, by
\Cref{thm:orthogonality}, not relatedness-faithful: related processes
receive maximally distant digests, and the relatedness is unrecoverable
from $\dig$ alone.
\end{corollary}

\begin{proof}
By construction $\dig$ realises avalanche, the extreme of
identity-faithfulness; \Cref{thm:orthogonality} forbids it from being
relatedness-faithful.
\end{proof}

\subsection{Time-blindness of a content hash}

The deeper limitation of a content hash for scheduling is not only that it
destroys structural relatedness, but that it ignores \emph{when} a process
ran.

\begin{definition}[Time-blind labelling]
\label{def:timeblind}
A labelling $\lab$ is \emph{time-blind} if it is a function of process
content alone: for any two processes with identical content but distinct
trajectory states, $\lab$ assigns the same label.
\end{definition}

\begin{theorem}[Content hashes are time-blind]
\label{thm:timeblind}
A cryptographic content hash $\dig$ is time-blind: it takes no temporal
input, so two processes with identical content and distinct trajectory
counts $\Mcount_1\neq\Mcount_2$ satisfy $\dig(p_1)=\dig(p_2)$.
\end{theorem}

\begin{proof}
$\dig$ is by definition a function of the content bytes only
\citep{nist180-4,katz2014}; the trajectory count is not among its inputs.
Two processes with identical content present identical inputs to $\dig$,
so $\dig(p_1)=\dig(p_2)$ regardless of $\Mcount_1,\Mcount_2$. The
labelling is therefore constant on content, i.e.\ time-blind by
\Cref{def:timeblind}.
\end{proof}

\begin{corollary}[Temporal relatedness is invisible to a content hash]
\label{cor:no-temporal}
Two processes close in trajectory count (occurring in the same burst of
work) but of unrelated content receive maximally distant digests; two
processes of identical content but far apart in time receive identical
digests. Neither temporal proximity nor temporal distance is legible from
the digest.
\end{corollary}

\begin{proof}
First claim: unrelated content means avalanche-distant digests
(\Cref{cor:no-related}); the small trajectory-count gap is not seen by
$\dig$ (\Cref{thm:timeblind}). Second claim: identical content gives an
identical digest (\Cref{thm:timeblind}) whatever the count gap.
\end{proof}

\subsection{The trajectory address supplies the missing axis}

The monotone count of \Cref{sec:count} is precisely the temporal input a
content hash lacks. We give it a prefix-structured form so that temporal
and structural proximity become prefix agreement.

\begin{definition}[Trajectory address]
\label{def:address}
Fix a branching factor $b\ge 2$. The \emph{trajectory address} of a
process at trajectory count $\Mcount$ along a structured work-tree is the
base-$b$ digit string $\Addr(\Mcount)=(a_1 a_2\cdots a_m)$, where the
digits are the process's choices at successive structural levels of the
tree (the path from the tree root to the process's current node). Two
addresses share a prefix of length $j$ iff the two processes agree on
their first $j$ structural choices.
\end{definition}

\begin{theorem}[Prefix is proximity]
\label{thm:prefix}
Let $\lcp(\Addr_1,\Addr_2)$ be the length of the longest common prefix of
two trajectory addresses. Then:
\begin{enumerate}[label=\textup{(\roman*)},leftmargin=2.2em,itemsep=2pt]
\item processes in the same structural subtree to depth $j$ have
$\lcp\ge j$ (structural relatedness $=$ prefix length);
\item along a single process, the address is monotone---each committed
unit appends one digit---so addresses order the process's history and
never collide (\Cref{thm:monotone});
\item $\lcp$ is computable in $\mathcal{O}(\min(m_1,m_2))$, i.e.\
$\mathcal{O}(1)$ in the address length, by digit comparison.
\end{enumerate}
\end{theorem}

\begin{proof}
(i) By \Cref{def:address}, the first $j$ digits encode the first $j$
structural choices; agreement on the subtree to depth $j$ is exactly
agreement on those $j$ digits, so $\lcp\ge j$. (ii) Each committed unit
descends one tree level, appending one digit; the count is strictly
increasing (\Cref{thm:monotone}), so successive addresses extend, never
repeat, and order the history. (iii) Longest-common-prefix of two strings
is a left-to-right digit scan, linear in the shorter length and bounded by
the fixed address width.
\end{proof}

\begin{remark}[Address vs.\ wall-clock timestamp]
\label{rem:vs-timestamp}
One could append a wall-clock timestamp to a digest to add a temporal
axis. The trajectory address is strictly preferable: it is monotone and
collision-free by \Cref{thm:monotone} (a wall clock can reset, run
backward, or collide within a tick); and its prefix encodes
\emph{structural} position, not merely an instant, so $\lcp$ measures
shared work, which a timestamp cannot
(\Cref{thm:prefix}(i)). It is the logical time of
\citep{lamport1978,fidge1988,mattern1989} carrying structural locality.
\end{remark}

\subsection{The paired label}

\begin{construction}[Content--time process label]
\label{con:paired}
Annotate each process $p$ with the ordered pair
\[
\lab(p)=\big(\,\dig(\mathrm{content}(p)),\ \Addr(\Mcount(p))\,\big),
\]
the cryptographic digest of its content and its trajectory address. The
two components are kept separate (not mixed into one digest).
\end{construction}

\begin{theorem}[Decidable relatedness queries]
\label{thm:queries}
Under \Cref{con:paired}, for any two processes $p_1,p_2$ the following are
decidable in $\mathcal{O}(1)$ (in label width):
\begin{enumerate}[label=\textup{(\roman*)},leftmargin=2.2em,itemsep=2pt]
\item \emph{recurrence:} $\dig_1=\dig_2$ and $\Addr_1\neq\Addr_2$ ---
same computation, occurring again at a later trajectory count
(\Cref{cor:noreturn}: a re-occurrence, not the original);
\item \emph{co-temporality:} $\dig_1\neq\dig_2$ and
$\lcp(\Addr_1,\Addr_2)$ large --- unrelated content sharing a structural
neighbourhood (a burst);
\item \emph{refinement:} $\dig_1\neq\dig_2$ (small content change) and
$\lcp(\Addr_1,\Addr_2)$ large --- a successive structural step of related
work.
\end{enumerate}
\end{theorem}

\begin{proof}
Each query is a constant number of comparisons on the two label
components: digest equality is an $\mathcal{O}(1)$ word comparison;
$\lcp$ is the bounded digit scan of \Cref{thm:prefix}(iii). (i) Digest
equality certifies identical content (collision-resistance of $\dig$
\citep{nist180-4,stevens2017}); distinct addresses certify distinct
trajectory counts (\Cref{thm:prefix}(ii)), hence a later occurrence. (ii)
Distinct digests certify unrelated content (\Cref{cor:no-related}); large
$\lcp$ certifies a shared subtree (\Cref{thm:prefix}(i)), i.e.\ temporal/%
structural co-location. (iii) Same as (ii) with the content difference
small, read as a refinement step rather than unrelated work.
\end{proof}

\begin{theorem}[Necessity of the pairing]
\label{thm:necessity}
No single component of \Cref{con:paired} decides all three queries.
\Cref{thm:queries}(i) (recurrence) is undecidable from $\Addr$ alone;
\Cref{thm:queries}(ii)--(iii) (co-temporality, refinement) are
undecidable from $\dig$ alone.
\end{theorem}

\begin{proof}
From $\Addr$ alone: two processes of \emph{identical} content but
committed at different counts have distinct addresses and are
indistinguishable, by the address, from two of \emph{different} content at
those counts---the address does not encode content, so recurrence
(same-content-again) cannot be certified. From $\dig$ alone: by
\Cref{cor:no-temporal} temporal proximity and refinement are invisible to
the content hash (time-blindness, \Cref{thm:timeblind}; relatedness
destruction, \Cref{cor:no-related}). Hence each axis is decided only by
its own component, and both components are required---the pairing is
necessary, matching \Cref{thm:orthogonality}.
\end{proof}

\begin{remark}[Two faces of a process]
\label{rem:two-faces}
The pairing reflects that a process is content-at-a-time: the digest
labels \emph{what} it is (maximal distinction, so it is unmistakably
itself---the identity face), the address labels \emph{when and where} in
the work-tree it sits (shared prefix with its kin---the relatedness
face). Identity wants maximal apartness; relatedness wants shared
structure; \Cref{thm:orthogonality} proves these cannot be one label, and
\Cref{con:paired} supplies both. If integrity over the address stream is
also required (tamper-evidence on the history), one additionally chains
the addresses in a Merkle hash \citep{merkle1988,damgaard1989}---a third,
separate use of a digest, for integrity, not relatedness.
\end{remark}

% =====================================================================
\section{Overhead and complexity of the scheduler}
\label{sec:overhead}
% =====================================================================

We bound the cost of running the scheduler and annotation themselves, so
that the machinery does not dominate the work it schedules.

\begin{theorem}[Per-tick cost]
\label{thm:pertick}
With $L$ live tasks and tick budget $B$ (work units), one call to
\textsc{Tick} (\Cref{alg:tick}) costs
$\mathcal{O}(B\log L)$ time and $\mathcal{O}(L)$ space, excluding the
dispatched work itself.
\end{theorem}

\begin{proof}
Priorities are maintained in a max-priority queue keyed by $P(\task)$;
each of at most $B$ dispatches per tick performs one $\argmax$
(queue-top, $\mathcal{O}(1)$) and one priority update on the dispatched
task ($\mathcal{O}(\log L)$ to re-heapify), giving $\mathcal{O}(B\log L)$.
The queue and task records are $\mathcal{O}(L)$ space. Reading the
returned residue and updating $\Delta$ over a fixed window is
$\mathcal{O}(1)$ per dispatch.
\end{proof}

\begin{theorem}[Annotation cost]
\label{thm:annot-cost}
Annotating one process (\Cref{con:paired}) costs one content-hash
evaluation plus $\mathcal{O}(1)$ to extend the trajectory address by one
digit. Each relatedness query (\Cref{thm:queries}) costs $\mathcal{O}(1)$
in label width.
\end{theorem}

\begin{proof}
The digest is one hash of the content; the address extends by one digit
per committed unit (\Cref{thm:prefix}(ii)). Queries are the bounded
comparisons of \Cref{thm:queries}.
\end{proof}

\begin{corollary}[Residue measurement does not dominate]
\label{cor:nondominate}
If reading a task's residue costs $\mathcal{O}(1)$ per unit (a single
functional evaluation, \Cref{def:residue}), the scheduler's overhead per
dispatched unit is $\mathcal{O}(\log L)$, independent of task running
times. The scheduler's cost is governed by the number of tasks, not by
how long any task takes---consistent with its refusal to predict
durations.
\end{corollary}

\begin{proof}
Per dispatched unit: $\mathcal{O}(\log L)$ for the queue update
(\Cref{thm:pertick}), $\mathcal{O}(1)$ for residue read and address
extension (\Cref{thm:annot-cost}). No term depends on the dispatched
work's duration.
\end{proof}

% =====================================================================
\section{Discussion and related work}
\label{sec:related}
% =====================================================================

\paragraph{Scheduling.} Classical scheduling minimises functions of
completion times from given processing times \citep{pinedo2016,
brucker2007,liu1973,coffman1968}. The residue scheduler removes the
given-processing-time assumption, which it argues is unattainable
(\Cref{sec:intro}), and replaces ranking-by-predicted-duration with
ranking-by-measured-residue and its descent
(\Cref{def:priority}). The closest prior stance is anytime computation
\citep{dean1988,zilberstein1996}, which trades solution quality for time;
we sharpen its stopping rule from ``time elapsed'' to ``residue stalled or
sufficient'' (\Cref{thm:progress,thm:sufficiency}) and give the
quality coordinate (residue) a floor (\Cref{thm:floor}).

\paragraph{Complexity.} Categorical complexity
(\Cref{def:complexity}) counts distinguishable work-trajectories and
yields a logarithmic termination bound for content-bounded tasks
(\Cref{thm:termination}); it refines, with explicit non-tight relations,
time/space complexity \citep{hartmanis1965,arora2009,sipser2013,blum1967},
and makes no class-collapse claim (\Cref{prop:standard}). The
unknowable-duration premise is the scheduling face of the absence of a
general running-time oracle \citep{turing1936,blum1967}.

\paragraph{Hashing and locality.} The label orthogonality theorem
(\Cref{thm:orthogonality}) is a separation between cryptographic hashing
(avalanche \citep{shannon1949,webster1986,feistel1973}, optimal for
identity \citep{nist180-4,katz2014,stevens2017}) and locality-sensitive
hashing (proximity-preserving \citep{indyk1998,broder1997,charikar2002},
optimal for relatedness). We do not propose a hash that does both---we
prove none can---and instead pair a content digest with a monotone
prefix-structured trajectory address, the latter a structurally enriched
logical clock \citep{lamport1978,fidge1988,mattern1989}. Time-blindness
of a content hash (\Cref{thm:timeblind}) is, to our knowledge, not
usually stated as a theorem though it is implicit in the definition of a
content-addressed store \citep{merkle1980,merkle1988}.

\paragraph{Bounded agents.} The Residue Floor Theorem
(\Cref{thm:floor}) is a finite-capacity result in the spirit of bounded
rationality \citep{simon1955}: a finite representation cannot resolve a
solution to a point, and the irreducible gap is the margin of
recognition. The information bound (\Cref{thm:info}) makes this a sizing
law on the residue scale \citep{shannon1948,cover2006}.

% =====================================================================
\section{Conclusion}
\label{sec:conclusion}
% =====================================================================

We have built a scheduler that does not predict running times---which we
argued, and the standard theory confirms, cannot be done in
general---and instead orders work by a measured residue and its rate of
descent. The construction rests on three proved facts: a bounded agent
cannot drive residue to zero, so ``solved'' is a region of positive
thickness $\floor>0$ whose boundary is the margin of recognition
(\Cref{thm:floor}); the count of committed work is strictly monotone, so a
process's progress is a faithful non-decreasing coordinate and a revisited
state is never a return (\Cref{thm:monotone}); and the number of
distinguishable work-trajectories grows as $d(d+1)^{n-1}$, giving
content-bounded tasks a logarithmic termination bound and a static
complexity baseline (\Cref{thm:inflation,thm:termination}). The live
deviation of residue from this baseline diagnoses, without any duration
estimate, whether a task is compute-limited (feed it) or structure-limited
(starve it) (\Cref{thm:diagnosis}), and a sub-task is released as soon as
it is sufficient for the parent goal, even above its own floor, by a
judgment the sub-task itself cannot make
(\Cref{thm:sufficiency,thm:nolocal}). The scheduler's priority rule is
sound, livelock-free, and of $\mathcal{O}(\log L)$ overhead per unit
(\Cref{thm:sound,thm:progress,cor:nondominate}).

For annotation we proved that identity and relatedness demand opposite
locality and so cannot share a label (\Cref{thm:orthogonality}): a
cryptographic digest is optimal for identity and destroys relatedness, and
is moreover structurally time-blind (\Cref{thm:timeblind}). Pairing the
digest with a monotone, prefix-structured trajectory address supplies the
missing temporal/structural axis and makes recurrence, co-temporality, and
refinement decidable in $\mathcal{O}(1)$ (\Cref{thm:queries}), with the
pairing proved necessary (\Cref{thm:necessity}). A process is
content-at-a-time; its label must carry both faces, and no single hash can.

The throughline is one quantity: the residue---the positive-thickness gap
between a partial result and a recognised solution. It is at once the
scheduler's progress signal, the floor that makes recognition possible,
the diagnostic that separates compute-limits from structure-limits, and
the basis on which a process is annotated by what it is and when it ran.
A scheduler that measures the gap, rather than predicting the time, needs
neither an oracle it cannot have nor a label that cannot exist.

% =====================================================================
\bibliographystyle{plainnat}
\bibliography{references}
% =====================================================================

\end{document}
